\reviewer%

\begin{generalcomment}
	This paper proposes VeriBoost, 
	a general acceleration approach for fault-tolerance verification in wide-area networks. 
	By employing link pruning and link compression, 
	the proposed method reduces the failure scenario space and improves the efficiency of various network verifiers. 
	The paper demonstrates certain practical value; 
	however, several concerns remain regarding its novelty, applicability, and experimental evaluation. 
	Major revision is recommended. The specific comments are as follows:
\end{generalcomment}
\begin{revresponse}[Thank you for your feedback.]
	We have carefully addressed all the issues item by item as follows.
\end{revresponse}

\begin{revcomment}
	The proposed method mainly relies on existing graph-theoretic techniques for link pruning and compression. 
	The authors should further clarify its key innovations and theoretical contributions compared with existing network verification acceleration methods.
\end{revcomment}
\begin{revresponse}
	Yes. When using \href{https://www.ctan.org/pkg/biblatex}{biblatex}, you can use the \texttt{refsection=section} option to achieve that.
	If we cite a new reference like \cite{Besser2021} here, it will again be number [1].

	Note that you might have to run \texttt{pdflatex} and \texttt{biber} multiple times.

	And reference [1] for \autoref{rev:2}~\cite{ReviewerReference} is now number [2].
	% \printpartbibliography{Besser2021,ReviewerReference}
\end{revresponse}

\begin{revcomment}
	The link compression method relies on specific topology structures and routing protocol assumptions. 
	The authors should further discuss the applicability of these assumptions in complex real-world WANs and their impact on the effectiveness of the proposed method.
\end{revcomment}
\begin{revresponse}
	Yes. When using \href{https://www.ctan.org/pkg/biblatex}{biblatex}, you can use the \texttt{refsection=section} option to achieve that.
	If we cite a new reference like \cite{Besser2021} here, it will again be number [1].

	Note that you might have to run \texttt{pdflatex} and \texttt{biber} multiple times.

	And reference [1] for \autoref{rev:2}~\cite{ReviewerReference} is now number [2].
	% \printpartbibliography{Besser2021,ReviewerReference}
\end{revresponse}

\begin{revcomment}
	Fault-tolerance verification is defined to cover scenarios with up to k link failures, 
	while some experimental analyses only consider combinations with exactly k failed links. 
	The authors should clarify whether all relevant failure scenarios are fully covered.
\end{revcomment}
\begin{revresponse}
	Yes. When using \href{https://www.ctan.org/pkg/biblatex}{biblatex}, you can use the \texttt{refsection=section} option to achieve that.
	If we cite a new reference like \cite{Besser2021} here, it will again be number [1].

	Note that you might have to run \texttt{pdflatex} and \texttt{biber} multiple times.

	And reference [1] for \autoref{rev:2}~\cite{ReviewerReference} is now number [2].
	% \printpartbibliography{Besser2021,ReviewerReference}
\end{revresponse}

\begin{revcomment}
	The experiments mainly use real network topologies with synthetic routing configurations. 
	The lack of evaluation using real-world complex routing configurations limits the evidence supporting the practical applicability of the proposed method.
\end{revcomment}
\begin{revresponse}
	Yes. When using \href{https://www.ctan.org/pkg/biblatex}{biblatex}, you can use the \texttt{refsection=section} option to achieve that.
	If we cite a new reference like \cite{Besser2021} here, it will again be number [1].

	Note that you might have to run \texttt{pdflatex} and \texttt{biber} multiple times.

	And reference [1] for \autoref{rev:2}~\cite{ReviewerReference} is now number [2].
	% \printpartbibliography{Besser2021,ReviewerReference}
\end{revresponse}

\begin{revcomment}
	The paper emphasizes the maximum speedup achieved, but provides insufficient performance analysis for networks with low link reduction ratios. A more comprehensive evaluation is needed to assess the actual acceleration benefits under different network conditions.
\end{revcomment}
\begin{revresponse}
	Yes. When using \href{https://www.ctan.org/pkg/biblatex}{biblatex}, you can use the \texttt{refsection=section} option to achieve that.
	If we cite a new reference like \cite{Besser2021} here, it will again be number [1].

	Note that you might have to run \texttt{pdflatex} and \texttt{biber} multiple times.

	And reference [1] for \autoref{rev:2}~\cite{ReviewerReference} is now number [2].
	% \printpartbibliography{Besser2021,ReviewerReference}
\end{revresponse}

\begin{changes}[III-A][4 left]
Yes. When using \href{https://www.ctan.org/pkg/biblatex}{biblatex}, you can use the \texttt{refsection=section} option to achieve that.
	If we cite a new reference like \cite{Besser2021} here, it will again be number [1].

	Note that you might have to run \texttt{pdflatex} and \texttt{biber} multiple times.

	And reference [1] for \autoref{rev:2}~\cite{ReviewerReference} is now number [2].

\end{changes}